\section{Application: Nickelate Superconductors}
  \label{sec:nickelates}

  \subsection{Structural Differences Relative to Cuprates}
    \label{subsec:nickelates-structure}

    Infinite-layer nickelates $R_{1-x}$Sr$_x$NiO$_2$
    share the same square lattice symmetry ($D_{4h}$)
    as cuprates,
    but differ in several structural respects.

    First,
    the antiferromagnetic structure factor
    at wavevector $(\pi,\pi)$
    is significantly reduced relative to cuprates.
    RIXS and neutron measurements indicate
    a smaller magnetic exchange scale
    $J_{\mathrm{Ni}}$
    and a weaker peak intensity
    $S_{\mathrm{Ni}}(\pi,\pi)$.

    Second,
    the Fermi surface is multi-band.
    In addition to the Ni-derived band,
    an electron pocket near $\Gamma=(0,0)$
    originating from rare-earth orbitals
    modifies the nesting properties
    and partially isotropizes the electronic structure.

    Within the framework of Section~\ref{sec:symmetry-selection-families},
    these differences imply a reduced
    frustration ratio
    \begin{equation}
      r_{\mathcal F}^{\mathrm{Ni}}
      <
      r_{\mathcal F}^{\mathrm{Cu}},
    \end{equation}
    where $r_{\mathcal F}$
    is taken as a first-order proxy
    proportional to $S(\pi,\pi)$.
    A quantitative definition of this proxy
    and representative experimental values
    are given in Appendix~\ref{app:frustration}.

  \subsection{Predicted Gap Symmetry}
    \label{subsec:nickelates-symmetry}

    The symmetry channel is determined
    by minimizing the cost functional
    $\mathcal{C}_{\Pi}$
    under the reduced and partially isotropized
    frustration constraint.

    In cuprates,
    the strong staggered $(\pi,\pi)$ frustration
    selects the $B_{1g}$ representation
    ($d_{x^2-y^2}$ symmetry).
    In nickelates,
    the reduced frustration amplitude
    and the presence of a $\Gamma$ pocket
    alter the minimization landscape.

    The framework predicts that
    the lowest-cost configuration at optimal doping
    lies in an extended $s$-wave class,
    specifically an $s^\pm$ structure
    with the following characteristics:

    \begin{enumerate}
      \item
      The gap is nodeless within each Fermi pocket
      at optimal doping.

      \item
      A sign change occurs between
      the Ni-derived and rare-earth-derived pockets.

      \item
      In the underdoped regime,
      nodal structure may develop near pocket edges
      as $r_{\mathcal F}$ increases
      and the staggered constraint regains angular weight.
    \end{enumerate}

    Many spin-fluctuation approaches yield
    either $d$-wave or multi-band $s^\pm$
    depending on assumed Fermiology
    and interaction strength,
    with no current consensus.
    The present mechanism differs in that
    symmetry selection is controlled
    by the real-space frustration ratio
    rather than by a momentum-dependent interaction kernel.

  \subsection{Quantitative Estimate of the Critical Temperature}
    \label{subsec:nickelates-tc}

    To leading order,
    the transition temperature in the quasi-2D regime
    obeys
    \begin{equation}
      k_B T_c
      \simeq
      \frac{\pi}{2}
      \delta J
      \, f(r_{\mathcal F}),
    \end{equation}
    as introduced in Eq.~\ref{eq:Tc-general}.

    For nickelates,
    using representative experimental values
    \begin{equation}
      J_{\mathrm{Ni}}
      \approx
      60\text{--}80 \,\mathrm{meV},
      \qquad
      \delta_{\mathrm{opt}}^{\mathrm{Ni}}
      \approx
      0.20,
    \end{equation}
    and taking
    $r_{\mathcal F}^{\mathrm{Ni}}$
    as a first-order proxy
    for the reduced $(\pi,\pi)$ structure factor,
    one obtains
    \begin{equation}
      \frac{T_c^{\mathrm{Ni}}}{T_c^{\mathrm{Cu}}}
      \approx
      \frac{\delta_{\mathrm{Ni}} J_{\mathrm{Ni}}}
      {\delta_{\mathrm{Cu}} J_{\mathrm{Cu}}}
      \cdot
      \frac{
        r_{\mathcal F}^{\mathrm{Ni}}
      }{
        r_{\mathcal F}^{\mathrm{Cu}}
      },
    \end{equation}
    up to order-unity corrections.

    Using representative cuprate values
    $J_{\mathrm{Cu}}\approx120\text{--}150\,\mathrm{meV}$
    and $\delta_{\mathrm{opt}}^{\mathrm{Cu}}\approx0.16$,
    this estimate yields
    \begin{equation}
      T_c^{\mathrm{Ni}}
      \sim
      0.3\text{--}0.4 \,
      T_c^{\mathrm{Cu}},
    \end{equation}
    consistent with observed nickelate
    transition temperatures in the 10--40\,K range,
    and potentially higher under pressure.

    This estimate involves no free fitting parameter.
    All quantities entering the ratio
    ($J$, $\delta$, $S(\pi,\pi)$)
    are independently measurable.

  \subsection{Disorder and Pressure Signatures}
    \label{subsec:nickelates-signatures}

    The predicted $s^\pm$ symmetry
    implies intermediate sensitivity
    to non-magnetic disorder.

    Specifically,
    the suppression rate
    \begin{equation}
      \frac{dT_c}{dn_{\mathrm{imp}}}
    \end{equation}
    under non-magnetic impurity substitution
    should lie between
    the Abrikosov--Gor'kov $d$-wave limit
    and the Anderson-theorem $s^{++}$ limit.
    Systematic impurity studies therefore
    provide a direct test of the mechanism.

    Under pressure,
    both the effective exchange scale $J_{\mathrm{Ni}}$
    and the inter-orbital hybridization
    are modified,
    entering the framework
    through $J$ and $r_{\mathcal F}$.
    An enhancement of $T_c$
    correlated with measurable changes
    in magnetic excitation spectra
    would support the structural interpretation.

  \subsection{Experimental Falsifiability}
    \label{subsec:nickelates-falsifiability}

    The mechanism would be invalidated if:

    \begin{itemize}
      \item
      A robust $d_{x^2-y^2}$ gap symmetry
      is conclusively established in nickelates
      despite reduced $(\pi,\pi)$ frustration.

      \item
      The measured suppression of $T_c$
      under non-magnetic disorder
      matches precisely either
      pure $d$-wave
      or fully Anderson-protected $s^{++}$ behavior.

      \item
      $T_c$ enhancement under pressure
      occurs without measurable modification
      of magnetic excitation scales
      or frustration proxies.
    \end{itemize}

    The predictive content of the framework
    lies in the enforced coupling
    between frustration amplitude,
    symmetry selection,
    and phase stiffness scaling.
